\documentclass{article}
\usepackage[utf8]{inputenc}
\usepackage[spanish]{babel}
\usepackage{listings}
\usepackage{xcolor}
\usepackage{geometry}
\geometry{a4paper, margin=1in}

% Configuración para código C++
\definecolor{codegreen}{rgb}{0,0.6,0}        
\definecolor{codegray}{rgb}{0.5,0.5,0.5}     
\definecolor{codepurple}{rgb}{0.58,0,0.82}   
\definecolor{backcolour}{rgb}{0.95,0.95,0.92}

\lstset{
    backgroundcolor=\color{backcolour},      
    commentstyle=\color{codegreen},
    keywordstyle=\color{magenta},
    numberstyle=\tiny\color{codegray},       
    stringstyle=\color{codepurple},
    basicstyle=\ttfamily\small,
    breakatwhitespace=false,
    breaklines=true,
    captionpos=b,
    keepspaces=true,
    numbers=left,
    numbersep=5pt,
    showspaces=false,
    showstringspaces=false,
    showtabs=false,
    tabsize=2,
    language=C++,
    inputencoding=utf8,
    extendedchars=true,
    literate={á}{{\'a}}1 {é}{{\'e}}1 {í}{{\'i}}1 {ó}{{\'o}}1 {ú}{{\'u}}1 {ñ}{{\~n}}1 {¿}{{\textquestiondown}}1 {¡}{{\textexclamdown}}1
}

% Macros para las secciones de los apuntes
\newcommand{\quees}[1]{
    \section*{🤔 ¿QUÉ ES?}
    \hrulefill \\
    #1
}

\newcommand{\dichodeotraforma}[1]{
    \section*{💬 DICHO DE OTRA FORMA}
    \hrulefill \\
    #1
}

\newcommand{\diagrama}[1]{
    \section*{🎨 DIAGRAMA}
    \hrulefill \\
    #1
}

\newcommand{\comofunciona}[1]{
    \section*{⚙️ ¿CÓMO FUNCIONA?}
    \hrulefill \\
    #1
}

\newcommand{\complejidad}[1]{
    \section*{📱 COMPLEJIDAD (Big-O)}
    \hrulefill \\
    #1
}

\newcommand{\entrevistas}[1]{
    \section*{🎯 EN ENTREVISTAS}
    \hrulefill \\
    #1
}

\newcommand{\resumen}[1]{
    \section*{📋 RESUMEN RÁPIDO}
    \hrulefill \\
    #1
}

\begin{document}

\title{Queues (Filas / Colas)}
\author{Aldo Zepeda Montoya}
\date{\today}
\maketitle

\subsection*{CATEGORÍA: 01_Estructuras_Lineales}

\quees{
La fila para entrar al cine 🎟️ o al banco 🏦. El primero que llega 
es el primero en ser atendido. No puedes "colarte" al frente; llegas 
al final y esperas tu turno.

Es una estructura "FIFO" (First In, First Out): el primero que entra
es el primero en salir. Es lo opuesto al Stack.
}

\dichodeotraforma{
Una Queue es una estructura lineal que restringe el acceso:
- Solo puedes AGREGAR (Enqueue) por el final (Back/Rear).
- Solo puedes QUITAR (Dequeue) por el frente (Front).

Existen variaciones importantes:
1. Deque: Puedes insertar/eliminar por ambos extremos.
2. Priority Queue: Los elementos salen según su importancia, no 
   solo por orden de llegada (implementada usualmente con un Heap).
}

\diagrama{\begin{verbatim}
enqueue(A) → enqueue(B) → enqueue(C)

FRONT                    BACK
  ↓                       ↓
 [A] ──→ [B] ──→ [C]

dequeue() → devuelve A y lo quita, queda:

FRONT         BACK
  ↓             ↓
 [B] ──→ [C]
\end{verbatim}}

\begin{lstlisting}[language=C++]
1. enqueue(x): Agregas el elemento `x` al final de la fila. O(1).
2. dequeue(): Quitas el primer elemento del frente. O(1).
3. front() / peek(): Miras quién es el primero sin quitarlo. O(1).
4. isEmpty(): Verificas si hay alguien en la fila. O(1).
5. size(): Devuelves cuántas personas hay esperando. O(1).

📝 PSEUDOCÓDIGO
──────────────────────────────────────────────────────────────

// BFS (Breadth First Search) básico usando una Queue
función BFS(inicio):
    fila = crear Queue
    fila.enqueue(inicio)
    marcar inicio como visitado

    MIENTRAS fila NO esté vacía:
        actual = fila.dequeue()
        PARA cada vecino de actual:
            SI vecino NO visitado:
                fila.enqueue(vecino)
                marcar vecino como visitado

💻 CÓDIGO C++
──────────────────────────────────────────────────────────────

#include <iostream>
#include <queue>    // STL Queue
#include <deque>    // STL Deque (Double-ended queue)
using namespace std;

int main() {
    // 1. Queue estándar (FIFO)
    queue<int> q;
    
    q.push(10); // Enqueue
    q.push(20);
    q.push(30);

    cout << "Primero en la fila: " << q.front() << endl; // 10
    cout << "Último en la fila: " << q.back() << endl;   // 30

    q.pop(); // Dequeue (quita el 10)
    cout << "Nuevo primero: " << q.front() << endl;     // 20

    // 2. Deque (Doble extremo)
    deque<int> dq;
    dq.push_back(100);
    dq.push_front(50); // [50, 100]
    dq.pop_back();     // [50]
    
    cout << "Frente de deque: " << dq.front() << endl;

    return 0;
}

⏱️ COMPLEJIDAD (Big-O)
──────────────────────────────────────────────────────────────

  Operación        │ Tiempo  │ Espacio
  ─────────────────┼─────────┼──────────────
  enqueue(x)       │ O(1)    │ O(1)
  dequeue()        │ O(1)    │ O(1)
  front() / peek() │ O(1)    │ O(1)
  isEmpty()        │ O(1)    │ O(1)
  Total (n datos)  │ -       │ O(n)
\end{lstlisting}

\entrevistas{
✅ Sin Queue no hay BFS: Si el problema pide "camino más corto" en un
   grafo sin pesos o explorar un árbol por niveles, vas a usar una Queue.

💡 Priority Queue: Úsala cuando necesites el "máximo" o "mínimo" 
   elemento de forma constante mientras insertas nuevos datos.

⚠️ Implementación con Array: Cuidado, si implementas una Queue con un
   Array simple, el `dequeue` podría ser O(n) si tienes que mover todos 
   los elementos. Por eso se usan "Circular Queues".

🔑 "Sliding Window Maximum": Un problema clásico de Google que se 
   resuelve de forma óptima usando un Deque.
}

\resumen{
• FIFO: First In, First Out.
• Operaciones O(1) en el frente y el final.
• Ideal para: Procesamiento de tareas, buffers, algoritmos de grafos
  como BFS.
• Deque: Permite meter/sacar por ambos lados (muy flexible).
• Priority Queue: Saca al que tiene más "rango", no al que llegó primero.
════════════════════════════════════════════════════════════════
}

\end{document}